%=============================================================================
% Section 8: The Evolution of AI Systems
%=============================================================================
\section{The Evolution of AI Systems}

\subsection{Historical Progression}

The introduction of the contact layer reflects a broader evolution in AI system design:

\textbf{Phase 1: Model Intelligence (2017--2022)}

Focus: Improving model capabilities

Examples: GPT-2, GPT-3, BERT, T5

Characteristics: Models as standalone capabilities

\textbf{Phase 2: Capability Integration (2022--2025)}

Focus: Connecting models to tools and services

Examples: ChatGPT plugins, LangChain, AutoGPT

Characteristics: Models with tool use

\textbf{Phase 3: Environment Interaction}

Focus: AI systems interacting with environments

Examples: AI-native applications, autonomous systems, edge AI

Characteristics: Action layers, capability ecosystems

Status: Current and emerging phase

This evolution can be summarized as: Intelligence $\to$ Capabilities $\to$ Contact Layer $\to$ Action.

\subsection{Future Trajectory}

As AI systems mature, we anticipate:
\begin{enumerate}[label=\arabic*.]
\item \textbf{Standardized Capability Protocols:} Industry-wide protocols for capability description
\item \textbf{Capability Marketplaces:} Economic ecosystems for capability provision and consumption
\item \textbf{Edge-First Architecture:} Contact layers deployed at the edge for latency-sensitive applications
\item \textbf{Physical-Digital Integration:} Seamless interaction with both digital services and physical devices
\end{enumerate}

The contact layer concept provides the architectural foundation for these developments.
